% =====================================================================
\chapter{Arborescence du projet}
\label{ch:arbo}

Ce chapitre parcourt l'arborescence réelle. Le format demandé est respecté :
\textbf{dossier/fichier \arrowto{} rôle \arrowto{} dépendances \arrowto{}
utilisation}.

\section{Vue générale}

\begin{lstlisting}[caption={Arborescence réelle --- les fichiers vides sont signalés}]
DXC-OASIS-Intelligent-IT-Operations-Assistant/
|-- docker-compose.yml        postgres:16 + chromadb  (SEULEMENT ces deux-la)
|-- README.md                 vue "Architecture Freeze V1" (anterieure au code actuel)
|-- .env                      POSTGRES_*, DATABASE_URL, CHROMA_*, JWT_*, GROQ_*, ANTHROPIC_*
|-- .claude/launch.json       lancement backend (8000) / frontend (5173)
|-- nginx/nginx.conf          [VIDE - 0 octet]
|-- .github/workflows/ci.yml  [VIDE - 0 octet]
|
|-- backend/
|   |-- Dockerfile            [VIDE - 0 octet]
|   |-- requirements.txt      dependances (dont mcp==2.0.0, paramiko, groq, chromadb)
|   |-- alembic.ini           configuration des migrations
|   |-- alembic/
|   |   |-- env.py
|   |   `-- versions/         14 migrations, chaine lineaire
|   `-- app/
|       |-- main.py           FastAPI, 11 routeurs, lifespan (3 taches de fond)
|       |-- api/
|       |   |-- deps.py       get_db, get_current_user, require_role
|       |   `-- v1/           11 modules de routes
|       |-- core/
|       |   |-- config.py       Settings pydantic-settings (lit .env)
|       |   |-- security.py     bcrypt + JWT HS256 + token cross-domain
|       |   |-- logging.py
|       |   |-- permissions.py  [VIDE]
|       |   `-- constants.py    [VIDE]
|       |-- database/
|       |   |-- base.py       declarative_base()
|       |   |-- session.py    engine + SessionLocal + get_db
|       |   |-- seed.py       3 roles + 3 comptes de test
|       |   `-- init_db.py    [VIDE]
|       |-- models/           16 classes SQLAlchemy sur 11 fichiers (+ 2 vides)
|       |-- schemas/          schemas Pydantic (plusieurs vides)
|       |-- services/         16 services metier
|       |-- personas/
|       |   |-- base_persona.py       [MORT : importe par personne]
|       |   `-- storage/
|       |       |-- agent.py          StoragePersona (445 lignes) : coeur de l'IA
|       |       |-- prompts.py        STORAGE_SYSTEM_PROMPT
|       |       `-- incident_rules.py seuils de detection
|       |-- mcp/
|       |   `-- storage/
|       |       |-- server.py       TOOL_REGISTRY + DESTRUCTIVE/REMEDIATION_TOOLS
|       |       |-- tool_server.py  MCPServer("oasis-storage-mcp"), transport stdio
|       |       |-- client.py       StorageMCPClient + audit mcp_calls
|       |       |-- ssh.py          paramiko : run(), run_many(), quote()
|       |       `-- tools/          capacity, backup, snapshot, restore,
|       |                           repair, health, performance
|       |-- rag/
|       |   |-- embedding_service.py  E5 multilingue, prefixes query:/passage:
|       |   |-- vectorstore.py        Chroma + garde-fou EmbeddingModelMismatch
|       |   |-- chunker.py            500 mots, overlap 50
|       |   |-- retriever.py          point d'entree unique des personas
|       |   |-- query_pipeline.py     assemblage du contexte RAG
|       |   |-- indexing_pipeline.py  indexation d'un dossier
|       |   |-- reindex.py            reconstruction complete des collections
|       |   |-- storage_kb/           4 documents markdown (NetApp, Veeam, TrueNAS, AWS)
|       |   `-- storage_runbooks/     4 runbooks correspondants
|       |-- dashboard/kpi_service.py  agregation du dashboard (461 lignes)
|       |-- orchestrator/
|       |   |-- llm_router.py           [MORT : Claude non branche]
|       |   |-- orchestrator.py         [VIDE]
|       |   `-- cross_domain_manager.py [VIDE]
|       |-- middleware/
|       |   |-- audit_middleware.py   ACTIF : journalise toute requete
|       |   |-- rbac_middleware.py    [MORT : doublon de deps.require_role]
|       |   `-- auth_middleware.py    [VIDE]
|       |-- llm/groq_client.py        client Groq + tool calling + fallback modele
|       `-- reports/generator.py      [VIDE - la vraie generation est dans services/]
|
|-- frontend/
|   |-- Dockerfile            [VIDE - 0 octet]
|   |-- vite.config.ts        alias @ -> src, proxy /api -> 127.0.0.1:8000
|   |-- tailwind.config.ts
|   `-- src/
|       |-- App.tsx           routes /storage/*, RequireAuth
|       |-- auth/AuthContext.tsx   session, signIn, signOut
|       |-- lib/
|       |   |-- api.ts         fetch + Bearer + download blob + upload FormData
|       |   |-- types.ts       259 lignes de types miroirs des schemas Pydantic
|       |   |-- useResource.ts hook de chargement + reload
|       |   |-- format.ts      dates, humanize
|       |   `-- cn.ts          clsx + tailwind-merge
|       |-- components/
|       |   |-- ActionProposal.tsx  boutons Confirm / Reject
|       |   |-- layout/    AppLayout, Sidebar, Topbar, GlobalSearch, PageHeader
|       |   |-- dashboard/ KpiCard, CapacityOverview, PerformanceChart,
|       |   |              RecentActivity, OasisInsights, AttentionBanner
|       |   `-- ui/        Badge, Button, Card, Field, Overlay, States, Toast
|       `-- pages/         Login, Dashboard, Chat, Incidents, Requests,
|                          DataHub, Reports, Settings
|
`-- vm-lab/
    |-- 00-network.sh         IP fixe host-only + sshd
    |-- 10-storage-lab.sh     LVM, LUN, restic, timers systemd
    |-- 20-nas-object.sh      NFS, Samba, MinIO, Object Lock
    |-- inject-incident.sh    provoque une panne pour la demonstration
    `-- README.md             correspondance "equipement d'entreprise" -> "objet Linux"
\end{lstlisting}

\clearpage

\section{Backend : le noyau}

\begin{longtable}{@{}p{3.2cm}p{4.1cm}p{3.5cm}p{3.4cm}@{}}
\toprule
\textbf{Fichier} & \textbf{Rôle} & \textbf{Dépendances} & \textbf{Quand}\\
\midrule
\endhead
\file{app/main.py} & Assemble l'application : middlewares, 11 routeurs, 3 tâches de fond & Tous les routeurs, \file{kpi\_service}, \file{incident\_watcher}, \file{metrics\_collector} & Au démarrage d'uvicorn\\
\file{core/config.py} & Objet \cls{Settings} unique lisant \file{.env} & \texttt{pydantic-settings} & Importé partout via \fn{settings}\\
\file{core/security.py} & \fn{hash\_password}, \fn{verify\_password}, \fn{create\_access\_token}, \fn{decode\_access\_token}, \fn{create\_cross\_domain\_token} & \texttt{passlib[bcrypt]}, \texttt{python-jose} & Connexion, chaque requête, audit\\
\file{api/deps.py} & \fn{get\_current\_user} et fabrique \fn{require\_role} & \file{security.py}, \file{session.py}, \cls{User} & Dépendance de presque toutes les routes\\
\file{database/session.py} & \fn{engine}, \fn{SessionLocal}, générateur \fn{get\_db} & SQLAlchemy & Chaque requête, chaque tâche de fond\\
\file{database/base.py} & \fn{Base = declarative\_base()} & SQLAlchemy & Hérité par les 16 classes de modèle\\
\file{database/seed.py} & Crée 3 rôles et 3 comptes de test & \cls{Role}, \cls{User}, passlib & Script manuel\\
\file{middleware/audit\_middleware.py} & Journalise chaque requête HTTP & \file{security.py}, \cls{AuditLog} & Toutes les requêtes sauf \file{/docs}, \file{/openapi.json}, \file{/redoc}, \file{/favicon.ico}\\
\bottomrule
\end{longtable}

\section{Backend : routes (\file{app/api/v1/})}

\begin{longtable}{@{}p{2.7cm}p{3.1cm}p{4.5cm}p{4.0cm}@{}}
\toprule
\textbf{Fichier} & \textbf{Préfixe} & \textbf{Rôle} & \textbf{Appelle}\\
\midrule
\endhead
\file{auth.py} & \route{/auth} & Connexion, profil courant & \fn{auth\_service.login}, \fn{get\_current\_user}\\
\file{storage.py} & \route{/storage} & Chat : CRUD des conversations + tour de dialogue & \fn{chat\_service}, \fn{storage\_persona}, \fn{incident\_investigation}\\
\file{actions.py} & \route{/storage/actions} & Propose / confirme / rejette une action & \fn{incident\_service}, \cls{StorageMCPClient}, \fn{incident\_verification}\\
\file{incidents.py} & \route{/incidents} & Liste, fiche détaillée, investigation, statut & \fn{incident\_service}, \fn{incident\_investigation}\\
\file{dashboard.py} & \route{/storage/dashboard} & Overview, performance, activité, insights & \fn{kpi\_service}\\
\file{knowledge.py} & \route{/storage/documents} & Data Hub : dépôt, lecture, suppression & \fn{knowledge\_service}\\
\file{reports.py} & \route{/storage/reports} & PDF global et PDF d'incident & \fn{report\_service}\\
\file{requests.py} & \route{/requests} & Demandes ingénieur / admin + pièces jointes & \fn{request\_service}\\
\file{search.py} & \route{/storage/search} & Recherche globale multi-entités & \fn{search\_service}\\
\file{admin\_routes.py} & \route{/admin} & Comptes, rôles, journaux d'audit (lecture) & \fn{admin\_service}, \fn{audit\_service}\\
\file{cross\_domain.py} & \route{/cross-domain} & Demandes d'accès temporaire & \fn{access\_request\_service}\\
\bottomrule
\end{longtable}

\section{Backend : services (\file{app/services/})}

\begin{longtable}{@{}p{3.5cm}p{1.0cm}p{10.5cm}@{}}
\toprule
\textbf{Fichier} & \textbf{L.} & \textbf{Rôle et moment d'utilisation}\\
\midrule
\endhead
\file{incident\_service.py} & 367 & Cœur transactionnel : \fn{create\_action\_request}, \fn{confirm\_and\_execute\_action}, \fn{reject\_action\_request}, \fn{create\_ticket\_from\_diagnosis}, \fn{resolve\_ticket}, \fn{add\_note}, \fn{has\_completed\_action}. Appelé par l'agent, les routes actions et incidents, et le watcher.\\
\file{incident\_investigation.py} & 254 & Ouvre le fil d'investigation, fait tourner l'agent avec le contexte du ticket, synthétise \texttt{root\_cause} et \texttt{impact\_summary} par un second appel Groq, propose une action de repli.\\
\file{incident\_verification.py} & 151 & Relit sur la VM l'indicateur qui avait ouvert le ticket. Renvoie un \cls{Verdict} (\texttt{resolved}, \texttt{detail}, \texttt{conclusive}).\\
\file{incident\_watcher.py} & 264 & Tâche de fond (300\,s) : détecte les jobs en échec, ouvre les tickets, referme ceux dont la remédiation confirmée a produit son effet.\\
\file{chat\_service.py} & 184 & Persistance des fils : \fn{save\_turn}, \fn{history\_for\_llm} (10 derniers messages), \fn{attach\_actions}, \fn{title\_from\_message}.\\
\file{request\_service.py} & 494 & Demandes : direction déduite du rôle, table \fn{ALLOWED\_TRANSITIONS}, pièces jointes stockées sous UUID.\\
\file{report\_service.py} & 489 & Génération PDF avec ReportLab : \fn{build\_global\_report}, \fn{build\_incident\_report}.\\
\file{knowledge\_service.py} & 235 & Data Hub : validation, extraction de texte, découpage, embeddings, \texttt{upsert} Chroma, statut d'indexation.\\
\file{search\_service.py} & 233 & Recherche ILIKE échappée sur 4 entités, avec extrait contextuel.\\
\file{metrics\_service.py} & 144 & \fn{collect\_once}, \fn{purge\_old}, \fn{read\_series} (agrégation \texttt{date\_bin} en deux passes).\\
\file{metrics\_collector.py} & 64 & Boucle de fond (300\,s) et purge quotidienne (rétention 31 jours).\\
\file{mcp\_audit.py} & 70 & Point d'écriture unique de \tbl{mcp\_calls}.\\
\file{access\_request\_service.py} & 80 & Workflow cross-domain : création, approbation avec jeton temporaire, refus.\\
\file{auth\_service.py} & 21 & \fn{authenticate\_user} et \fn{login}.\\
\file{admin\_service.py} & 15 & Lecture des comptes et des rôles.\\
\file{audit\_service.py} & 16 & Lecture des journaux d'audit.\\
\bottomrule
\end{longtable}

\section{Frontend}

\begin{longtable}{@{}p{3.4cm}p{4.2cm}p{3.4cm}p{3.2cm}@{}}
\toprule
\textbf{Fichier} & \textbf{Rôle} & \textbf{Dépendances} & \textbf{Quand}\\
\midrule
\endhead
\file{main.tsx} & Monte \cls{App} dans \texttt{\#root} & React DOM & Chargement\\
\file{App.tsx} & Routeur + \cls{AuthProvider} + \cls{ToastProvider} + \cls{RequireAuth} & react-router-dom & Toujours\\
\file{auth/AuthContext.tsx} & Session : \fn{signIn}, \fn{signOut}, reprise via \route{/auth/me} & \file{lib/api.ts} & Au montage et à la connexion\\
\file{lib/api.ts} & Client HTTP unique & \texttt{fetch} & Chaque appel réseau\\
\file{lib/useResource.ts} & Hook \texttt{data / error / loading / reload} & \cls{ApiError} & Dashboard, incidents, demandes, Data Hub\\
\file{components/ActionProposal.tsx} & Confirmation humaine & \file{api.ts}, \cls{Badge}, \cls{Button}, \cls{Toast} & Chat et fiche d'incident\\
\file{components/layout/GlobalSearch.tsx} & Barre de recherche et navigation & \route{/storage/search} & Topbar\\
\file{pages/ChatPage.tsx} & 738 lignes : liste de fils, envoi, renommage, suppression & \file{api.ts}, \cls{ActionProposal} & Route \route{/storage/chat}\\
\file{pages/IncidentsPage.tsx} & 521 lignes : filtres, fiche, cycle de vie, investigation & \file{api.ts}, \cls{ActionProposal} & Route \route{/storage/incidents}\\
\file{pages/RequestsPage.tsx} & 1318 lignes : boîtes « sent » et « received », pièces jointes & \file{api.ts} & Route \route{/storage/requests}\\
\bottomrule
\end{longtable}

\section{Laboratoire \file{vm-lab/}}

Ce dossier n'est pas de la documentation : ce sont les scripts qui
\emph{fabriquent} l'infrastructure que les outils interrogent. Le principe est
énoncé dans son README : « l'application ne simule rien \ldots{} c'est donc la
machine qu'il faut équiper --- pas le code ».

\begin{longtable}{@{}p{4.6cm}p{10.2cm}@{}}
\toprule
\textbf{Équipement d'entreprise} & \textbf{Équivalent construit dans la VM}\\
\midrule
\endhead
Disques d'une baie SAN & Fichiers loopback creux (\texttt{truncate} + \texttt{losetup}) déclarés en PV\\
Pool de la baie & Volume group \texttt{vg\_prod}\\
LUN présenté à un hyperviseur & Logical volume formaté, monté sous \file{/srv/luns/}\\
Volume NFS d'un NAS & \texttt{nfs-kernel-server}, export \file{/srv/nas/nfs\_prod}\\
Partage SMB & Samba, partage \texttt{finance} authentifié\\
Stockage objet S3 & MinIO\\
Logiciel de sauvegarde (Veeam) & \texttt{restic} piloté par des timers systemd\\
Cible cloud hors-site & Dépôt restic sur MinIO\\
Sauvegarde immuable & Object Lock MinIO en mode COMPLIANCE\\
Snapshot de baie et restauration & \texttt{lvcreate -{}-snapshot}, \texttt{lvconvert -{}-merge}\\
\bottomrule
\end{longtable}

\begin{infobox}
Le README du laboratoire liste aussi honnêtement \textbf{ce qui n'est pas
simulable} : Fibre Channel et zonage (aucune émulation FC n'existe), robots de
bandes LTO (\texttt{mhvtl} exigerait la compilation d'un module noyau) et
microcodes de baie. C'est exactement le genre de limite qu'un jury apprécie de
voir assumée.
\end{infobox}
